\scanpage{58}
% archive foliation: 57
3: Relations avec courbes elliptiques:
\inserted{\uline{le morphisme} $f_2^t:\Sigma_{\underline J}\longrightarrow\bar\Sigma_{\underline J}$.}

On suppose tjrs les car. rés. de $S\ne2$. Soient $E_0$ une courbe elliptique sur $S$, $E$ un torseur sous $E_0$ (de sorte que $E_0\simeq\underline{\operatorname{Pic}}^0_{E/S}$).

On suppose donné dans $E$ un $S$-autom
\[
\tag{17}\sigma:E\longrightarrow E.
\]
tel que
\[
\tag{18}\sigma_{E_0}^{\unsure{u|h}}=-\operatorname{id}_{E_0}
\]
ce qui implique que loc. (ét.), identifiant $E$ à $E_0$ via une section origine de $E$, on aura
\[
\sigma x=a-x\qquad a\text{ section convenable}
\]
ce qui est d'ailleurs intrinsèque si on interprète $a$ comme une section de $E\wedge_{E_0}E\simeq\underline{\operatorname{Pic}}^2_{E/S}$.

On aura donc
\[
\sigma^2=\operatorname{id}_E.
\]
Soit \deleted{$\unsure{\phi|I}$}
\[
\tag{19}\left\{\begin{array}{ll}
\underline I=E^\sigma&\text{sous-schéma des pts fixes de }\sigma\\
\underline I\subset E&\\[-4pt]
\Vert&\\[-4pt]
\{x\in E\mid 2x=a\}&
\end{array}\right.
\]
on voit que $\underline I$ est un torseur sous le schéma en groupes
\[
\tag{20}\underline V={}_2E_0=\operatorname{Ker}(2\operatorname{id}_{E_0}:E_0-E_0),
\]
qui est justement un groupe \inserted{sur $S$} loc. isom. à $\mathbb D_{2S}$.

La donnée de $(E,\sigma)$, où $E$ est un schéma
